\section{The Architect's Question}\label{the-architects-question}

Designing an AI system means producing candidate architectures from the
substrate --- memory, compute, communication, and parallelism --- rather
than selecting a canned stack. This chapter works through a synthesis
method: given a characterized workload from the six-dimension framework
(Ch. 4), generate two or three candidate architectures that each satisfy
the hard constraints (memory fit, latency budget, economics), then carry
them forward for evaluation in Chs. 14--16. The core thesis is that
design is constraint-driven generation, not menu-picking. This
distinguishes our approach from the abstract meta-method presented in
Ch. 22, which operates at a level of generality that does not
instantiate against concrete substrate constraints. In the language of
the book's spine --- the decision loop \emph{constrain → synthesize →
measure → decide → refine} --- this chapter is the \textbf{synthesize}
step: we turn a constraint-characterized workload into a bounded set of
candidate architectures that then enter the measure/decide stages in
Chs. 14--16.

\section{Concept}\label{concept}

The architect's first step is to reverse-engineer the workload into its
constituent constraints. From the six dimensions --- quality
requirements, traffic profile, token profile, latency budget, economic
ceiling, and operational constraints --- we extract the three hard
constraints that any candidate must satisfy: (1) the total weight memory
must fit the target residency budget, (2) the KV-cache footprint plus
weight memory must fit the host's aggregate memory, and (3) the total
request latency (prefill + decode) must stay within the budget. Any
architecture that violates any one of these is eliminated before
evaluation. This constraint-first posture inverts the typical practice
of browsing vendor offerings and then hoping they fit; instead, we
generate architectures that we know can physically reside and operate
within the specified substrate.

\section{Mental Model}\label{mental-model}

Think of the substrate as a set of bounded resources and the
architecture design problem as a packing exercise. The resources are:
GPU memory (weight residency + KV cache), memory bandwidth (prefill
throughput), inter-GPU communication volume (for parallelism schemes),
and the latency budget split between prefill and decode. A candidate
architecture is a valid packing of the workload's weight memory and
KV-cache into the available GPU memory, paired with a parallelism and
caching strategy that delivers the required tokens-per-second throughput
within the latency budget. The architect alternates between squeezing
the residency footprint (e.g., via KV quantization or precision
reduction) and adjusting the parallelism layout (e.g., parameter
vs.~data vs.~pipeline) until all three hard constraints are
simultaneously satisfiable. This mental model --- substrate as packing
problem --- is what enables systematic candidate generation rather than
ad hoc selection.

\section{Worked Example}\label{worked-example}

We generate three candidate architectures for the canonical
enterprise-Q\&A RAG workload: \textasciitilde10 requests/s average, peak
40 rps, 70B FP16 model, 9.2K token input (1,200 prompt + 8K retrieved
context), 300-token output. The canonical weight memory is 140 GB (70B ×
2 bytes FP16). The KV-cache for 9.2K context at FP16 precision ({[}1P
§14{]}) is approximately 24 GB (2 × num\_layers × hidden\_dim × 9,200 ×
2 bytes / scale factor), giving a total residency of 164 GB. A single
8×H100 host provides 640 GB aggregate memory (8 × 80 GB), leaving ample
headroom. Each candidate satisfies the three hard constraints
differently.

\textbf{Candidate (a): single 8×H100 host with prefix caching (Ch.
11).}\\
All 70B weights and the full 24 GB KV cache reside on the 8-GPU host.
Prefix caching from Ch. 11 amortizes the prefill attention over repeated
requests with the same prompt prefix, reducing effective prefill FLOPs
by \textasciitilde30\% at 10 rps average and \textasciitilde50\% at peak
40 rps when many requests share common document roots. The 8×H100 host
fits the full 164 GB residency easily, but on \emph{throughput} we must
be honest with the \hyperref[chap:8]{Chapter 8} arithmetic: a single node sustains only
\textasciitilde2,471 prefill-tokens/s, far below the \textasciitilde92K
input tokens/s the workload needs on average --- prefill is
compute-bound here, and meeting it is exactly the scaling problem
Chapters 15--17 take up. Memory fit and throughput fit are different
gates: the memory gate passes with one host; the throughput gate does
not. Memory fit: 140 GB weights + 24 GB KV = 164 GB \textless{} 640 GB
aggregate ({[}1P §14{]}). Latency budget: prefill + decode stay within
budget thanks to prefix caching and the high bandwidth of 8×H100.
Economics: one host purchase, one software license tier.

\textbf{Candidate (b): P/D-disaggregated two-pool.}\\
We keep the full 70B model in both pools and split by \emph{phase} ---
which is what prefill/decode (P/D) disaggregation actually means
(\hyperref[chap:11]{Chapter 11}): a compute-bound prefill pool and a bandwidth-bound decode
pool, each holding the complete 140 GB weights plus the in-flight KV,
with the KV hidden-state handed between them (Splitwise / DistServe /
Mooncake). Each pool therefore needs its own 8×H100 host (residency
\textasciitilde164 GB --- weights plus one request's KV --- comfortably
under 640 GB). The win is that prefill and decode no longer fight over
the same host's resources: the prefill pool runs FLOPS-saturated on the
9.2K input, the decode pool runs bandwidth-saturated on the 300 output
tokens, and KV crosses between them. The cost is a second host plus the
KV-transfer fabric and orchestration, so this pays only once the
workload actually outgrows a single host (\hyperref[chap:11]{Chapter 11}'s ``when to
disaggregate''). Economics: two hosts and added fabric; it is the right
architecture at scale, not for the single-host canonical case.

\textbf{Candidate (c): KV-quantized smaller host.}\\
We downgrade to a 7B-class model (7B × 2 bytes = 14 GB weights) with
KV-cache quantized to 4-bit, yielding a KV footprint of \textasciitilde6
GB for 9.2K context. Total residency: 14 GB + 6 GB = 20 GB, which fits
on a single H100 (80 GB) with 60 GB of headroom for overhead. The
trade-off is a significant quality drop: the 7B model cannot achieve the
same answer quality on complex enterprise Q\&A as the 70B model, even
with retrieval augmentation. To compensate, we apply KV-quantization at
4-bit, which recovers most of the lost quality relative to unquantized
KV, but the model's intrinsic capacity remains the binding constraint.
Throughput on a single H100 is far higher than on the 70B --- a 7B
prefill costs only \textasciitilde14 GFLOP/token, so one card sustains
tens of thousands of prefill tokens/s (vs \textasciitilde2,471 for the
70B node, \hyperref[chap:8]{Chapter 8}) --- yet even that cannot change the verdict: at
peak 40 rps the workload needs \textasciitilde368K input tokens/s, so
satisfying peak prefill still needs more than a handful of cards, and
the 7B's quality ceiling on complex enterprise Q\&A remains the binding
constraint regardless of throughput. This candidate illustrates the
extreme end of the KV-quantization space: the memory footprint is
trivially small, but neither quality nor throughput makes it competitive
for this workload.

\emph{Memory/residency arithmetic summary.} The canonical weights (140
GB) + KV (24 GB FP16) = 164 GB fits 640 GB (8×H100). Candidate (a) uses
the full host with room to spare. Candidate (b) keeps the full model in
each pool, so each 8×H100 host also holds \textasciitilde164 GB.
Candidate (c) fits 20 GB on one H100, but its quality ceiling rules it
out for the quality target.

\section{Measurement}\label{measurement}

To validate a candidate architecture, the architect must derive three
numbers from first principles and verify them against the substrate: 1.
\textbf{Weight memory.}

\[
W = N \times \text{bytes-per-param}
\]

where \(N\) is the parameter count and bytes-per-param is 2 for FP16
(replace 2 with bit-width/8 for quantization). Derived fact {[}2°
DERIVED{]}: \(70 \times 10^9 \times 2 \text{ B} = 140\) GB. 2.
\textbf{KV-cache footprint.}

\[
KV = 2 \times n_\text{layers} \times d_\text{hidden} \times L \times \text{bytes-per-activation}
\]

where \(L\) is the context length. For a 70B model with 8-bit KV
caching, this is approximately 12 GB for 9.2K context; at FP16 precision
it doubles to \textasciitilde24 GB {[}2° DERIVED{]}. The architect
should compute this for the exact context length of the workload, not
assume a fixed value. 3. \textbf{Throughput per device.}
Tokens-per-second that a single device can sustain for prefill and
decode, measured on the target hardware. This is the number that,
multiplied by the number of devices in the candidate layout, must exceed
the workload's tokens-per-second demand (92K input tokens/s average, 3K
output tokens/s average, 368K input tokens/s peak, 12K output tokens/s
peak). Vendors publish peak numbers; the architect should treat these as
{[}VERIFY{]} benchmarks and measure on equivalent hardware when
possible.

All three numbers must be confirmed before a candidate is carried
forward ({[}1P §14{]}). If any one fails, the candidate is returned to
the synthesis step.

\section{Common Mistakes}\label{common-mistakes}

\begin{itemize}
\tightlist
\item
  \textbf{Menu-picking without constraint validation.} Browsing vendor
  offerings and selecting the one that ``feels right'' rather than
  generating candidates against quantified hard constraints. This is the
  opposite of the constraint-driven method described here.
\item
  \textbf{Ignoring the KV-cache residency arithmetic.} Assuming that the
  weight memory is the only memory budget item. In input-heavy workloads
  ({[}1P §14{]}: 9.2K average context, 30× more input than output), the
  KV cache can be 15--20\% of total residency and is often the binding
  constraint.
\item
  \textbf{Over-optimizing prefix caching.} Assuming that prefix caching
  can amortize prefill across all requests. In practice, cache
  effectiveness depends on the overlap of prompt prefixes across the
  traffic workload; sparse prefix sharing reduces the realized savings.
\item
  \textbf{Using KV-quantization as a free lunch.} Quantizing the KV
  cache reduces memory footprint but introduces quality degradation that
  may be unacceptable for the target workload. The architect must
  measure the quality impact before committing to a quantized KV layout.
\item
  \textbf{Forgetting the latency split.} Treating the latency budget as
  a single number rather than a prefill/decode split. A candidate may
  meet the total latency budget by trading prefill for decode or vice
  versa, which may not satisfy the service-level objectives for either
  leg.
\end{itemize}

\section{Architecture Consequence}\label{architecture-consequence}

The three candidates carry forward into the evaluation chapters (Chs.
14--16) with distinct test plans. Candidate (a), the single 8×H100 host
with prefix caching, is the baseline: we evaluate throughput,
cost-per-token, and latency SLO compliance on a single hardware
configuration. Candidate (b), the P/D-disaggregated two-pool, is
evaluated on the added complexity of inter-group communication and the
pipeline overlap efficacy; we measure whether the latency budget is
maintained under the partitioned layout. Candidate (c), the KV-quantized
7B model, is evaluated on the quality-degradation trade-off; we compare
answer quality against the 70B baseline and compute the economic
break-even point (number of hosts required) against the quality loss.
The synthesis method produces these candidates as a set from which the
evaluation chapter selects based on quantitative results, not on gut
feeling.

\section{What We Still Don't Know}\label{what-we-still-dont-know}

As of 2026-08, several questions remain open and are flagged
{[}VERIFY{]}: - \textbf{KV-quantization quality impact at 4-bit for
enterprise Q\&A.} The quality loss from 4-bit KV quantization has not
been systematically characterized across the diversity of enterprise
prompts in the canonical workload. {[}VERIFY HYPOTHESIS{]} -
\textbf{Prefix-caching effectiveness under realistic traffic overlap.}
The assumed 30--50\% prefill amortization rests on the premise that many
requests share document-level prefixes. Empirical measurement of prefix
overlap across a production enterprise traffic trace is needed to
confirm these numbers. {[}VERIFY HYPOTHESIS{]} -
\textbf{P/D-disaggregation communication overhead in multi-host
settings.} The boundary between parameterparallel and data-parallel
placement is relatively unexplored at the 70B scale; the actual
communication volume and its latency impact depend on the routing
library and network topology. {[}VERIFY HYPOTHESIS{]} - \textbf{Whether
8-bit KV caching is sufficient for 9.2K context at FP16-equivalent
quality.} The arithmetic in §3 assumes 8-bit KV reduces memory by 2×
with minimal quality loss, but the interaction between 8-bit
quantization and the retrieval-augmented generation pipeline has not
been verified. {[}VERIFY HYPOTHESIS{]}

These flags exist because home-lab measurements are not reference per
the evidence taxonomy; they will be resolved (promoted, dropped, or
demoted) before publication.

\section{End-of-Chapter Mini-Case}\label{end-of-chapter-mini-case}

An architect is brought into a design conversation for a company-wide
internal Q\&A system. The stakeholder states: ``We need an AI system
that can answer employee questions over our internal documentation, with
acceptable quality and within budget.'' Before any architecture can be
defended, the architect applies the constraint-driven method from this
chapter.

First, the workload is characterized across the six dimensions. Quality
requirements fix the model class: the system must achieve
\textasciitilde70B-model answer quality on enterprise prompts. The
traffic profile is estimated at \textasciitilde500 employees active
during business hours, generating \textasciitilde20 questions/hour peak,
\textasciitilde2 rps average. The token profile is 1,000-token prompts
plus 8K retrieved context plus 300-token answers, roughly 9.3K tokens
per request. The latency budget is 5 seconds end-to-end. The economic
ceiling is \$15,000 monthly operating budget. The operational constraint
is that the system must run on existing on-premise GPU infrastructure
(no cloud add-ons).

From the token profile, the architect computes: 9.3K input tokens + 300
output tokens ≈ 9.6K tokens per request. At 2 rps average and 20 rps
peak, the system must sustain \textasciitilde19.2K input tokens/s
average and \textasciitilde384K input tokens/s peak. The hard
constraints are then extracted: weight memory must fit the on-premise
GPU inventory; KV cache + weights must fit within that inventory; total
tokens-per-second across the installed devices must exceed the peak
traffic demand.

The architect generates three candidates using the method from §3.
Candidate (a): a single 8×H100 host with 70B FP16 weights + 24 GB KV
cache and prefix caching, which fits the memory budget (164 GB
\textless{} 640 GB aggregate ({[}1P §14{]})) --- but as \hyperref[chap:8]{Chapter 8} shows,
one node sustains only \textasciitilde2,471 prefill-tokens/s, far short
of the \textasciitilde19.2K input tokens/s the mini-case's
\textasciitilde2 rps average demands, so the memory gate passes while
the throughput gate already fails. Candidate (b): a P/D-disaggregated
two-pool keeping the full 70B in each pool (prefill and decode pools,
each \textasciitilde164 GB residency on its own 8×H100 host), which buys
phase-splitting at the cost of a second host plus the KV-transfer
fabric. Candidate (c): a KV-quantized 7B model on a single H100, which
fits the memory budget (20 GB total residency) but fails the quality
constraint --- the 7B model cannot reach the required answer quality
regardless of KV quantization --- and fails the throughput demand at
peak, making it economically non-viable.

The architect discards Candidate (c) immediately: it satisfies the hard
memory constraint but violates the quality and economic constraints.
Between Candidates (a) and (b), the decision hinges on whether the
doubled hardware cost of Candidate (b) is justified by its operational
flexibility (e.g., independent scaling of prefill and decode pools,
fault isolation). The architect presents both candidates, together with
the honest throughput arithmetic of \hyperref[chap:8]{Chapter 8} (no canonical single host
meets the full prefill demand alone), to the stakeholder as the decision
foundation.

\begin{center}\rule{0.5\linewidth}{0.5pt}\end{center}

\begin{figure}
\centering
\pandocbounded{\includegraphics[keepaspectratio,alt={\hyperref[fig:12.1]{Fig 12.1} --- Candidate architecture synthesis: three candidate architectures from the canonical enterprise-Q\&A RAG workload, comparing memory residency, parallelism plan, and constraint-satisfaction outcome (memory fit, latency SLO, economics) {[}ILLUSTRATIVE conceptual{]}},width=\maxwidth{\textwidth},keepaspectratio]{figures/fig-12-1201.pdf}}
\caption{\hyperref[fig:12.1]{Fig 12.1} --- Candidate architecture synthesis: three candidate
architectures from the canonical enterprise-Q\&A RAG workload, comparing
memory residency, parallelism plan, and constraint-satisfaction outcome
(memory fit, latency SLO, economics) {[}ILLUSTRATIVE conceptual{]}}

\label{fig:12.1}\end{figure}

\subsection{\hyperref[tab:12.1]{Table 12-1} --- Candidate architectures with
constraint-satisfaction
columns}\label{table-12-1-candidate-architectures-with-constraint-satisfaction-columns}

\begin{longtable}[]{@{}
  >{\raggedright\arraybackslash}p{0.1667\linewidth}
  >{\raggedright\arraybackslash}p{0.1667\linewidth}
  >{\raggedright\arraybackslash}p{0.1667\linewidth}
  >{\raggedright\arraybackslash}p{0.1667\linewidth}
  >{\raggedright\arraybackslash}p{0.1667\linewidth}
  >{\raggedright\arraybackslash}p{0.1667\linewidth}@{}}
\toprule\noalign{}
\begin{minipage}[b]{\linewidth}\raggedright
candidate
\end{minipage} & \begin{minipage}[b]{\linewidth}\raggedright
memory (weights + KV)
\end{minipage} & \begin{minipage}[b]{\linewidth}\raggedright
host layout
\end{minipage} & \begin{minipage}[b]{\linewidth}\raggedright
throughput
\end{minipage} & \begin{minipage}[b]{\linewidth}\raggedright
SLO
\end{minipage} & \begin{minipage}[b]{\linewidth}\raggedright
economics
\end{minipage} \\
\midrule\noalign{}
\endhead
\bottomrule\noalign{}
\endlastfoot
(a) 8×H100 + prefix caching & 140 + 24 = \textbf{164 GB} & 1 × 8×H100 &
prefill \textasciitilde1,200 · decode \textasciitilde300 tok/s & ✓
memory 164 \textless{} 640 GB · ✓ latency in 5 s & ✓ 1 host \\
(b) P/D-disaggregated & 140 + 24 = \textbf{164 GB} per pool & 2 × 8×H100
& prefill pool \textasciitilde2,471 · decode BW-bound & ✓ memory per
host · ✓ phase-split at scale & ✗ 2 hosts + fabric \\
(c) KV-quantized 7B & 14 + 6 = \textbf{20 GB} & 1 × H100 &
\textasciitilde tens of thousands · decode \textgreater{} 70B & ✗ fails
quality ceiling & ✗ \textgreater1 card at peak \\

\label{tab:12.1}\end{longtable}

\emph{All figures trace to the §14 canonical scenario; the KV-cache and
throughput numbers are {[}2° DERIVED{]} worked examples, not measurement
claims.}

\begin{center}\rule{0.5\linewidth}{0.5pt}\end{center}

\emph{Next chapter: \hyperref[chap:13]{Chapter 13} --- Evaluating Candidate Architectures,
which subjects the three candidates from \hyperref[tab:12.1]{Table 12-1} to quantitative
throughput, cost, and quality evaluation across the canonical workload.}
